% GOVERNANCE NOTE::LSPV5-WRITE-ESCALATION-REVFIX1::2026-10-19T06:16:00+09:00
\documentclass[11pt]{article}
\usepackage[margin=1in]{geometry}
\begin{document}
\section*{LSPV-5 canonical-write escalation}
The newest verified canonical MAIN LaTeX artifact observed in the repository is
\texttt{research/main/2026-10-19\_0155\_MAIN\_NAELIFTKW2.tex}.  A later
\texttt{research/main/2026-10-19\_0450\_MAIN\_RELATIVEGAUGEFOLD23\_BRIEF.md}
exists, but it is not the required canonical \texttt{.tex} snapshot and does not
carry the required MAIN header/final sentinel.  The current handoff reports two
blocked synthesis write attempts.  No exact missing synthesis body is available
to LSPV-1, so no mathematical content has been reconstructed.

\paragraph{Failure classification.}
\textbf{UNKNOWN / probable prompt--procedure or body-availability failure.}
There is no path collision at an intended 04:50 canonical \texttt{.tex} path,
and repository reads succeed.  The evidence is insufficient to call this a
permission failure or stale-SHA conflict.

\paragraph{Smallest repair.}
On its next run LSPV-5 should preserve any exact outbox body that exists, fetch
current repository state, write it once to a fresh canonical \texttt{.tex} path
with the required MAIN markers, and verify readback before root reconciliation.
If the exact body does not exist, it must not be recreated from the brief; the
next run should instead produce one new current synthesis report normally.  The
rollback is simply removal of this procedural note after a later verified MAIN
records the resolution.  Scientific targets and schedules are unchanged.

P versus NP remains open.
% FINAL_VERIFICATION_SENTINEL::GOVERNANCE::LSPV5-WRITE-ESCALATION-REVFIX1::P_VS_NP_OPEN
\end{document}
